\chapter{Proposed Methodology}
\label{ch:methods}

\begin{chapterabstract}
This chapter develops the formal machinery of the thesis. It defines a
composite semantic-state representation built from cheap per-frame signals, a
temporal semantic memory that stabilises that representation against noise, a
bounded semantic-drift estimator that fuses embedding, object, and tracking
cues, a three-tier adaptive compute-allocation policy driven by interpretable
thresholds, and the Semantic Compute Efficiency metric used to evaluate the
resulting trade-off.
\end{chapterabstract}

\section{Overview and Design Principles}
\label{sec:meth-overview}

The proposed system follows a \emph{monitor-then-allocate} discipline. Every
frame passes through a lightweight semantic monitor that updates an internal
estimate of the scene's semantic state and emits a scalar drift score. A
scheduling policy maps this score to a discrete compute tier, and only the
upper tiers invoke the expensive vision-language model. Figure~\ref{fig:arch}
shows the end-to-end flow.

Three principles guide the design. First, \emph{the cost of deciding must be
much smaller than the cost of reasoning}; hence the monitor reuses
perception primitives that are cheap enough to run on every frame. Second,
\emph{decisions must be interpretable}; the policy is governed by two
thresholds with clear semantics rather than an opaque learned controller.
Third, \emph{the reasoning model is a black box}; the scheduler never inspects
VLM internals, so any captioning or visual-question-answering model can be
substituted without changing the policy.

\begin{figure}[htbp]
\centering
\begin{tikzpicture}[
  node distance = 6.5mm and 9mm,
  font = \footnotesize,
  box/.style = {rectangle, rounded corners=2pt, draw, thick,
                align=center, minimum height=7mm, inner sep=3pt,
                text width=30mm, fill=blue!5},
  tier1/.style = {box, fill=green!8},
  vlm/.style  = {box, fill=red!8, text width=34mm},
  arrow/.style = {-{Stealth}, thick},
]
\node[box] (cam) {Robot Camera /\\Video Stream};
\node[tier1, below=of cam] (filter) {Lightweight Frame\\Filtering (CPU)};
\node[tier1, below=of filter] (yolo) {YOLOv8-n Object\\Detection};
\node[tier1, below=of yolo] (track) {ByteTrack\\Object Tracking};
\node[tier1, below=of track] (clip) {CLIP Embedding\\$E_t$};

\node[box, right=22mm of track] (mem) {Temporal Semantic\\Memory $M_t$};
\node[box, above=of mem] (drift) {Semantic Drift\\Estimator $D_t$};
\node[box, above=of drift] (policy) {Adaptive Allocation\\Policy};
\node[vlm, below=of mem] (vlm) {Selective VLM\\Invocation (Tier 2/3)};
\node[box, below=of vlm] (cache) {Semantic Cache\\Update};
\node[box, below=of cache] (robot) {Robotic Perception\\Layer};

\draw[arrow] (cam) -- (filter);
\draw[arrow] (filter) -- (yolo);
\draw[arrow] (yolo) -- (track);
\draw[arrow] (track) -- (clip);
\draw[arrow] (clip.east) -- ++(0.6,0) |- (mem.west);
\draw[arrow] (yolo.east) to[out=0,in=180] (drift.west);
\draw[arrow] (mem) -- (drift);
\draw[arrow] (drift) -- (policy);
\draw[arrow] (policy.east) -- ++(0.5,0) |- (vlm.east);
\draw[arrow] (vlm) -- (cache);
\draw[arrow] (cache) -- (robot);
\draw[arrow] (cache.west) -- ++(-0.5,0) |- (mem.south west);
\end{tikzpicture}
\caption[System architecture]{End-to-end architecture. Green blocks form the
always-on Tier-1 monitor; the red block is the selectively invoked
vision-language model. The semantic cache closes the loop by feeding reused
outputs back into temporal memory.}
\label{fig:arch}
\end{figure}

\section{Semantic State Representation}
\label{sec:meth-state}

At each frame $t$ the system summarises the scene into a composite
\emph{semantic state}
\begin{equation}
  Z_t = f\!\left(E_t,\, O_t,\, M_t\right),
  \label{eq:state}
\end{equation}
where the three constituents are deliberately heterogeneous so that they
capture complementary notions of ``what the scene is.''

\begin{itemize}
  \item $E_t \in \mathbb{R}^{d}$ is the $\ell_2$-normalised CLIP image
  embedding~\cite{Radford2021CLIP} (for ViT-B/32, $d = 512$). It encodes
  global, appearance-and-semantics content in a metric space where cosine
  similarity is meaningful.
  \item $O_t = \{(c_i, n_i)\}$ is the multiset of object classes $c_i$ and
  their counts $n_i$ reported by the detector~\cite{Jocher2023YOLOv8}. It
  encodes \emph{symbolic} scene content --- what kinds of things are present
  and how many.
  \item $M_t$ is the temporal semantic memory (Section~\ref{sec:meth-memory}),
  a smoothed reference state that represents what the system currently
  ``believes'' the scene to be.
\end{itemize}

The state is intentionally a fusion of a continuous embedding and a symbolic
object set: embeddings are sensitive to holistic change but can be fooled by
illumination, whereas object sets are robust to appearance but blind to
configuration. Combining them yields a more faithful change signal than
either alone.

\section{Temporal Semantic Memory}
\label{sec:meth-memory}

A naive comparison of consecutive frames would react to every flicker of
sensor noise. To suppress this, the system maintains an exponentially decayed
memory of the embedding,
\begin{equation}
  M_t = \lambda\, M_{t-1} + (1-\lambda)\, E_t,
  \qquad 0 < \lambda < 1,
  \label{eq:memory}
\end{equation}
re-normalised to unit length after each update. The decay factor $\lambda$
sets the memory horizon: values near $1$ produce a long, stable memory that
ignores transient fluctuations, while smaller values track the scene more
responsively. Drift is measured against this stabilised memory rather than
against the immediately preceding frame, so that brief perturbations do not
trigger spurious reasoning while sustained change still accumulates. When the
VLM is invoked and the cache refreshed, the memory is anchored to the
freshly reasoned-over state, marking it as the new reference.

\section{Semantic Drift Estimation}
\label{sec:meth-drift}

The scalar drift signal fuses three bounded components,
\begin{equation}
  D_t = \alpha\, D_{\text{embed}} + \beta\, D_{\text{objects}}
        + \gamma\, D_{\text{track}},
  \qquad \alpha + \beta + \gamma = 1,
  \label{eq:drift}
\end{equation}
with each term normalised to $[0,1]$ so that $D_t \in [0,1]$ and the weights
form a convex combination.

\paragraph{Embedding drift.} The cosine distance between the current
embedding and the memory,
\begin{equation}
  D_{\text{embed}} = \tfrac{1}{2}\bigl(1 - \cos(E_t, M_{t-1})\bigr)
  = \tfrac{1}{2}\bigl(1 - E_t^{\!\top} M_{t-1}\bigr),
  \label{eq:dembed}
\end{equation}
which lies in $[0,1]$ for unit-norm vectors and grows as the holistic content
departs from the remembered state.

\paragraph{Object novelty.} Let $O_t$ and $O_{\text{ref}}$ be the current and
reference object multisets. Their dissimilarity is the Jaccard-style novelty
\begin{equation}
  D_{\text{objects}} = 1 - \frac{\lvert O_t \cap O_{\text{ref}} \rvert}
                                 {\lvert O_t \cup O_{\text{ref}} \rvert},
  \label{eq:dobjects}
\end{equation}
computed over class-count pairs so that both the appearance of a new class and
a change in the number of instances of an existing class register as novelty.

\paragraph{Tracking-identity change.} Using the identities maintained by
ByteTrack~\cite{Zhang2022ByteTrack}, let $\mathcal{I}_t$ be the set of active
track identities. The identity-churn score
\begin{equation}
  D_{\text{track}} = \frac{\lvert \mathcal{I}_t \,\triangle\, \mathcal{I}_{\text{ref}} \rvert}
                          {\lvert \mathcal{I}_t \cup \mathcal{I}_{\text{ref}} \rvert}
  \label{eq:dtrack}
\end{equation}
measures the symmetric difference of identities, capturing objects that have
entered or left the scene even when the class-level object set is unchanged
(for example, one person replacing another).

The three components are deliberately redundant in benign cases and
complementary in adversarial ones: appearance-only change is caught by
$D_{\text{embed}}$, count change by $D_{\text{objects}}$, and identity change
by $D_{\text{track}}$. The weights $(\alpha, \beta, \gamma)$ are treated as
hyperparameters and studied in the ablation of
Section~\ref{sec:exp-ablation}.

\section{Egocentric Operation and Camera-Motion Compensation}
\label{sec:meth-egomotion}

A robot perceives the world through an \emph{egocentric}, continuously moving
camera, so the moving-camera case is the primary deployment regime rather than
a corner case. Egomotion poses a specific threat to the drift signal: as the
camera pans, tilts, or translates, the global appearance --- and hence the
CLIP embedding $E_t$ --- changes even when the underlying scene is
semantically static. Left uncorrected, this inflates $D_{\text{embed}}$ and
produces \emph{false} transitions that would trigger needless reasoning,
eroding the very compute savings the system is designed to deliver.

Two design choices make the formulation egomotion-aware. First, two of the
three drift components are inherently robust to camera motion: the object-set
novelty $D_{\text{objects}}$ depends on \emph{what} is present, not on the
viewpoint, and the tracking-identity churn $D_{\text{track}}$ is preserved by
the multi-object tracker, which maintains identities across the smooth
viewpoint shifts that egomotion induces. The fragile component is the holistic
embedding term.

Second, the embedding term is explicitly \emph{discounted} by the fraction of
inter-frame change attributable to camera motion. A lightweight global-motion
estimator --- sparse optical flow or a homography fit between consecutive
downscaled frames --- yields a camera-motion fraction $c_t \in [0,1]$, where
$c_t \to 1$ when the change between frames is well explained by a single
global geometric transform (pure egomotion) and $c_t \to 0$ when a large flow
residual indicates genuine, non-global content change. The corrected embedding
drift is
\begin{equation}
  D_{\text{embed}}^{\text{ego}} = (1 - c_t)\, D_{\text{embed}},
  \label{eq:dembed-ego}
\end{equation}
which is substituted for $D_{\text{embed}}$ in the fused signal of
Equation~(\ref{eq:drift}) when operating on egocentric streams. Pure panning
($c_t \!\approx\! 1$) is thereby suppressed, while a new object appearing in
view ($c_t \!\approx\! 0$) still registers. The estimator runs on downscaled
frames within the Tier-1 budget, so the correction adds negligible cost to the
always-on monitor.

This component is part of the system design for robotic deployment; its
empirical validation on egocentric (head-mounted and mobile-robot) streams,
where static-camera assumptions break down, is a dedicated part of the
cross-domain evaluation plan of Section~\ref{sec:exp-crossdomain}.

\section{Adaptive Compute-Allocation Policy}
\label{sec:meth-policy}

The drift signal is mapped to one of three compute tiers through two
thresholds $0 < \tau_{\text{low}} < \tau_{\text{high}} < 1$:
\begin{equation}
  \text{Tier}(D_t) =
  \begin{cases}
    1 \;\; (\text{cache reuse, no VLM}) & D_t < \tau_{\text{low}},\\[2pt]
    2 \;\; (\text{lightweight VLM}) & \tau_{\text{low}} \le D_t < \tau_{\text{high}},\\[2pt]
    3 \;\; (\text{full VLM reasoning}) & D_t \ge \tau_{\text{high}}.
  \end{cases}
  \label{eq:policy}
\end{equation}
Tier~1 serves the cached output directly and costs nothing beyond the monitor.
Tier~2 issues a cheaper VLM query (for example, a short prompt or reduced
generation length) appropriate to a moderate change. Tier~3 performs full
reasoning and refreshes the cache and memory. Algorithm~\ref{alg:scheduler}
gives the complete per-frame procedure.

\begin{algorithm}[htbp]
\caption{Semantic-aware per-frame compute allocation}
\label{alg:scheduler}
\begin{algorithmic}[1]
\Require frame $x_t$; thresholds $\tau_{\text{low}},\tau_{\text{high}}$;
         weights $\alpha,\beta,\gamma$; decay $\lambda$
\State $O_t \gets \textsc{Detect}(x_t)$;\quad
       $\mathcal{I}_t \gets \textsc{Track}(O_t)$;\quad
       $E_t \gets \textsc{CLIP}(x_t)$
\State $D_{\text{embed}} \gets \tfrac{1}{2}(1 - E_t^{\!\top} M_{t-1})$
\State $D_{\text{objects}} \gets 1 - \text{Jaccard}(O_t, O_{\text{ref}})$
\State $D_{\text{track}} \gets \text{SymDiff}(\mathcal{I}_t, \mathcal{I}_{\text{ref}})$
\State $D_t \gets \alpha D_{\text{embed}} + \beta D_{\text{objects}} + \gamma D_{\text{track}}$
\If{$D_t < \tau_{\text{low}}$}
  \State $y_t \gets \textsc{CacheLookup}()$ \Comment{Tier 1: reuse}
\ElsIf{$D_t < \tau_{\text{high}}$}
  \State $y_t \gets \textsc{VLM}_{\text{light}}(x_t)$;\;
         \textsc{CacheUpdate}$(y_t)$ \Comment{Tier 2}
\Else
  \State $y_t \gets \textsc{VLM}_{\text{full}}(x_t)$;\;
         \textsc{CacheUpdate}$(y_t)$;\;
         $O_{\text{ref}},\mathcal{I}_{\text{ref}} \gets O_t,\mathcal{I}_t$ \Comment{Tier 3}
\EndIf
\State $M_t \gets \text{normalize}(\lambda M_{t-1} + (1-\lambda) E_t)$
\State \Return $y_t$
\end{algorithmic}
\end{algorithm}

\section{Selective Invocation and Semantic Caching}
\label{sec:meth-cache}

The semantic cache stores the most recent VLM output together with the state
that produced it. On a Tier-1 frame the cached output is returned unchanged,
preserving the information delivered to downstream tasks at zero reasoning
cost. On Tier-2 and Tier-3 frames the cache is refreshed and, in the Tier-3
case, the reference object set and track identities are re-anchored to the
current frame so that subsequent drift is measured relative to the
newly understood scene. This closes the feedback loop drawn in
Figure~\ref{fig:arch}.

\section{Semantic Compute Efficiency}
\label{sec:meth-sce}

Evaluating the scheduler requires a single figure of merit that rewards
saving compute only insofar as semantic fidelity is preserved. We define
\emph{Semantic Compute Efficiency} as
\begin{equation}
  \mathrm{SCE} = \frac{\text{SemanticRetention}}{\text{NormalizedComputeCost}},
  \label{eq:sce}
\end{equation}
where \emph{SemanticRetention} is the average similarity between the output
the system actually delivers (cached or freshly computed) and the output an
every-frame oracle would deliver, and \emph{NormalizedComputeCost} is the
total compute spent by the policy divided by that of the every-frame
baseline. A policy that reasons on every frame has retention $1$ and
normalised cost $1$, giving $\mathrm{SCE}=1$; a useful scheduler achieves
$\mathrm{SCE} > 1$ by retaining most of the fidelity at a fraction of the
cost. This metric is the primary axis of comparison in
Chapter~\ref{ch:experiments}.

\section{The Semantic Observatory and Policy-Discovery Methodology}
\label{sec:meth-observatory}

The thresholds and weights of the policy in Section~\ref{sec:meth-policy}
could be set by intuition, but a stronger scientific stance is to
\emph{discover} them from evidence about how semantic meaning actually evolves
in a stream. To that end this thesis introduces the \emph{Semantic
Observatory}: an instrument that consumes the per-frame log produced by a
pipeline run and derives a suite of semantic-evolution signals that go beyond
the instantaneous drift used by the scheduler. From the logged drift $D_t$ and
object/track metadata it computes, per frame, the quantities listed in
Table~\ref{tab:observatory-signals}, and it aggregates them into higher-level
structure: contiguous \emph{stability windows} (sustained low drift),
\emph{transition zones} (sustained high drift), per-track \emph{object
persistence} records, and discrete \emph{semantic events}.

\begin{table}[htbp]
\centering
\caption[Semantic Observatory signals]{Semantic-evolution signals derived by
the Semantic Observatory from a pipeline run.}
\label{tab:observatory-signals}
\small
\begin{tabular}{@{}ll@{}}
\toprule
\textbf{Signal} & \textbf{Definition} \\
\midrule
Semantic drift $D_t$        & Fused drift signal (Eq.~\ref{eq:drift}) \\
Acceleration $A_t$          & $D_t - D_{t-1}$ \\
Velocity                    & Smoothed rate of drift change \\
Volatility $V_t$            & Variance of $D_t$ over a rolling window \\
Scene entropy               & Shannon entropy of the object-class distribution \\
Object persistence          & Per-track lifespan, in frames \\
Stability windows           & Contiguous low-drift regions \\
Transition zones            & Contiguous high-drift regions \\
Semantic progress           & Cumulative semantic evolution \\
VLM trigger density         & Rolling VLM invocation rate \\
\bottomrule
\end{tabular}
\end{table}

These signals turn an opaque stream into an auditable record of \emph{why}
each scheduling decision was taken, and they expose the statistical structure
--- how long scenes stay stable, how abruptly transitions arrive, whether
spikes coincide with camera motion --- that a good policy should exploit. The
Observatory is deliberately \emph{read-only}: it observes and analyses but
never alters the scheduler, which prevents the circularity of tuning a policy
against a moving evaluation.

The Observatory underpins a disciplined experimental protocol for deriving
policies, summarised as a single loop:
\begin{equation}
  \textsf{Observe} \rightarrow \textsf{Hypothesise} \rightarrow
  \textsf{Modify one variable} \rightarrow \textsf{Evaluate} \rightarrow
  \textsf{Accept/Reject} \rightarrow \textsf{Record}.
  \label{eq:protocol}
\end{equation}
Each iteration begins from Observatory evidence, advances a single falsifiable
hypothesis (for example, ``camera motion inflates embedding drift''), changes
exactly one element of the monitor or policy, and is accepted only if it
improves SCE without degrading retention on held-out streams. Confining each
experiment to one variable keeps the source of any change attributable, and
recording every experiment --- observations, hypothesis, metrics, verdict ---
builds the evidence trail that distinguishes a discovered policy from a guessed
one. The scientific contribution of this thesis is therefore twofold: the
specific scheduler of Section~\ref{sec:meth-policy}, and the Observatory-driven
\emph{methodology} by which such schedulers can be derived and justified.
